\documentclass[11pt]{article}
\usepackage[margin=1in]{geometry}
\usepackage{amsmath,amssymb,booktabs,tabularx,array,graphicx,xcolor,hyperref}

\hypersetup{
  colorlinks=true,
  linkcolor=blue!55!black,
  citecolor=blue!55!black,
  urlcolor=blue!55!black
}
\newcolumntype{Y}{>{\raggedright\arraybackslash}X}
\newcommand{\microir}{\textsc{micro-ir}}
\newcommand{\cstate}{\mathcal{S}}

\title{Heterogeneous Cognitive Transformers:\\
Computational and Epistemic Coprocessors}
\author{Working Draft}
\date{August 2026}

\begin{document}
\maketitle

\begin{abstract}
Language-model systems commonly group calculators, formal solvers, databases,
and search engines under the single label of tool use.  We argue that this
grouping hides two different failure modes and therefore two different control
problems.  A \emph{competence deficit} occurs when the required information is
available but a specialized computation would be more reliable or efficient; a
\emph{knowledge deficit} occurs when reliable, fresh, or source-specific
evidence is absent.  We propose a position and an evidence-gated research
program in which a Transformer remains the learned semantic and generative
substrate while computational and epistemic coprocessors can be engaged by
generation-time \emph{semantic interrupts}.  An interrupt normalizes an
emerging need into typed micro-IR, routes it to a bounded engine, and records
the result as provenance-bearing cognitive micro-state before generation
continues.  This proposal does not claim to invent automatic tool use or active
retrieval; rather, it unifies them as two tracks under one runtime contract and
asks when automatic, persistent, typed coupling improves over explicit tool
calls, planner--executor agents, program-aided reasoning, and upfront or active
retrieval.  We specify the architecture, locate its novelty relative to prior
work, define component-level and end-to-end falsification criteria, and lay out
an interleaved sequence of experiments from strict reflexes to a heterogeneous
cognitive runtime.  Conventional tool calling and upfront retrieval remain
supported baselines and fallbacks throughout.
\end{abstract}

\section{Introduction}

Transformers are unusually capable semantic interfaces.  They map natural
language into useful representations, adapt to new tasks through context, and
generate coherent continuations.  They are nevertheless an awkward place to
perform some operations.  Exact arithmetic, logical closure, constraint
propagation, database aggregation, and source-sensitive factual verification
have algorithms and data structures whose reliability does not improve by
asking a language model to imitate them for more tokens.

The usual remedy is to give the model tools.  This is powerful and should not
be discarded.  However, the phrase \emph{tool use} covers systems with very
different timing, state, and control semantics: a model may explicitly emit an
API call, a planner may schedule a solver, a retriever may run once before
generation, or a learned token may request evidence during generation.  Treating
all of these as equivalent obscures a useful systems question:

\begin{quote}
Can specialized computation and evidence acquisition become ordinary,
observable parts of decoding, while the Transformer remains responsible for
semantic interpretation, composition, and communication?
\end{quote}

We call the resulting design a \emph{heterogeneous cognitive Transformer}.
The word \emph{cognitive} describes the role of the external state in the
system's ongoing inference; it is not a claim about consciousness or a
biological model.  The word \emph{coprocessor} emphasizes a typed, bounded,
runtime-managed subsystem that can be activated without requiring the language
model to author a complete natural-language tool-use plan.

This paper makes four contributions as a position paper, not as an empirical
result:
\begin{enumerate}
  \item It separates competence deficits from knowledge deficits and defines
  computational and epistemic coprocessors as twin experimental tracks.
  \item It specifies a common interrupt contract: detection, typed
  normalization, routing, execution, provenance-bearing state update, and
  controlled reinjection.
  \item It distinguishes the proposal from adjacent tool-use, retrieval, and
  neuro-symbolic systems without claiming that any individual mechanism is new.
  \item It turns the architecture into a falsifiable, interleaved research
  program with matched baselines, component metrics, scaling tests, and an
  evidence gate before every increase in coupling.
\end{enumerate}

\section{Two Deficits, Two Coprocessor Families}

\subsection{Competence deficits}

A competence deficit is present when the premises and operands are already
available, but a specialized procedure can perform an operation more reliably,
efficiently, or audibly than neural continuation.  Examples include arithmetic,
symbolic simplification, type checking, graph reachability, logical inference,
constraint solving, theorem proving, and bounded simulation.  The relevant
question is not ``what fact is missing?'' but ``what operation should be
performed over the available state?''

A \emph{computational coprocessor} consumes a typed operation and inputs and
returns a result, derivation, certificate, counterexample, or explicit failure.
Its answer is conditional on its inputs and semantics.  Formal derivability
does not establish that natural-language premises were translated correctly or
that those premises are true.

\subsection{Knowledge deficits}

A knowledge deficit is present when a pending claim requires information that
is missing, stale, source-specific, or insufficiently supported.  Examples
include current events, private database values, exact quotations, policy
versions, and claims for which provenance matters.  The relevant question is
``what evidence is needed, from which source, and with what freshness?''

An \emph{epistemic coprocessor} consumes a typed information need and returns
records, passages, graph neighborhoods, source metadata, timestamps, or an
explicit no-evidence result.  Retrieval supplies evidence, not truth.  Sources
can be stale, contradictory, irrelevant, or malicious; the runtime must retain
that uncertainty rather than silently converting retrieved text into fact.

\begin{table}[t]
\centering
\small
\caption{The distinction determines the request, success criterion, and main
failure mode.  Some tasks compose both families.}
\label{tab:deficits}
\begin{tabularx}{\linewidth}{@{}p{0.19\linewidth}YYY@{}}
\toprule
 & Competence deficit & Knowledge deficit & Composed case \\
\midrule
Need & Perform or check an operation & Acquire or verify evidence & Acquire operands, then compute \\
Typical request & Typed expression, constraints, graph, program & Entity, relation, query, source, freshness & Evidence request followed by typed operation \\
Primary output & Value, derivation, proof status, counterexample & Record, passage, provenance, evidence status & Provenance-linked derived result \\
Success test & Semantic and execution correctness & Retrieval relevance, support, freshness, source quality & Both tests plus correct composition \\
Characteristic risk & Wrong formalization or unsafe execution & Stale, conflicting, missing, or poisoned evidence & Laundering weak evidence through exact computation \\
\bottomrule
\end{tabularx}
\end{table}

The categories are not a forced partition of every request.  ``Which product
had the largest increase in sales last quarter?'' may require a database lookup
followed by aggregation.  The distinction remains useful because it exposes
where an error occurred and prevents a calculator result from being treated as
source evidence or a retrieved passage from being treated as a proof.

This symmetry is the program's broadest differentiator.  A competence deficit
routes toward calculator, logic, SMT, graph, algebra, or code; an epistemic
deficit routes toward relational DB, lexical index, vector store, knowledge
graph, or web search.  The shared question generalizes from retrieval timing to
which specialized computational or epistemic substrate can assist the evolving
neural computation now.

\subsection{Neural uncertainty is not epistemic risk}

Low next-token confidence is an observable model property; epistemic risk is a
property of a pending commitment's dependence on missing, fresh,
source-specific, or verifiable evidence.  Neither entails the other.  The
retrieval track must therefore measure confidence on the actual checkpoint and
populate four diagnostic regimes rather than assign them by intuition:
\begin{center}
\begin{tabular}{lll}
\toprule
Confidence & Epistemic risk & Diagnostic expectation \\
\midrule
High & Low & continue \\
Low & Low & possible confidence-triggered over-retrieval \\
High & High & critical semantic-trigger test \\
Low & High & both confidence and semantic triggers should fire \\
\bottomrule
\end{tabular}
\end{center}
Recall in the high-confidence, high-risk quadrant tests confident hallucination
and stale commitments. False retrieval in the low-confidence, low-risk quadrant
tests whether uncertainty is being mistaken for an evidence deficit.

\section{Architecture}

\subsection{The neural plane and the interrupt loop}

Let $y_{1:t}$ be the generated prefix and $h_t$ the decoding state available to
an interrupt policy.  The \emph{neural plane} is the Transformer-based semantic
and generative computation.  A watcher evaluates whether the current state
contains an actionable need:
\[
  d_t = D(y_{1:t}, h_t, \cstate_t), \qquad
  d_t \in \{\textsc{none},\textsc{compute},\textsc{retrieve}\}.
\]
If $d_t\neq\textsc{none}$, normalization produces a typed event
\[
  e_t = N(y_{1:t},h_t,\cstate_t)
      = (f,o,a,r,s,c,b),
\]
where $f$ is the family, $o$ the operation or information need, $a$ the typed
arguments, $r$ the epistemic role, $s$ the scope, $c$ a calibrated confidence,
and $b$ a resource budget.  A router selects an engine $k=R(e_t)$ or abstains.
The engine returns a typed outcome $z_t=E_k(e_t)$, and the state manager performs
an auditable update
\[
  \cstate_{t+1}=U(\cstate_t,e_t,k,z_t,\pi_t),
\]
where $\pi_t$ records provenance, versions, timing, and dependencies.  A
materializer $M$ exposes only the state needed for subsequent generation:
\[
  y_{t+1}\sim p_\theta\!\left(\cdot\mid y_{1:t},
  M(\cstate_{t+1},h_t)\right).
\]
Detection, normalization, routing, execution, update, materialization, and use
are separate failure sites and must be evaluated separately.

\begin{figure}[t]
\centering
\small
\fbox{\begin{minipage}{0.94\linewidth}
\centering
\textbf{Transformer neural plane: interpret $\rightarrow$ compose $\rightarrow$
generate}\\[4pt]
$y_{1:t},h_t$\\[-2pt]
$\Downarrow$\\[-2pt]
\textbf{Interrupt controller}\\
detect $\rightarrow$ normalize to typed \microir{} $\rightarrow$ route / abstain
\\[-2pt]
$\swarrow\hspace{0.31\linewidth}\searrow$\\[-2pt]
\begin{tabular}{p{0.40\linewidth}p{0.40\linewidth}}
\centering\textbf{Computational plane}\par
calculator, logic, constraints, graph, algebra, bounded code/simulation
&
\centering\textbf{Epistemic plane}\par
database, lexical/vector index, knowledge graph, web/search
\end{tabular}
\\[-2pt]
$\searrow\hspace{0.31\linewidth}\swarrow$\\[-2pt]
\textbf{Cognitive micro-state}\\
values, derivations, evidence, status, confidence, provenance, dependencies,
scope, version\\[-2pt]
$\Downarrow$\\[-2pt]
controlled textual or structured materialization $\rightarrow$ continued decoding
\end{minipage}}
\caption{The proposed runtime.  Explicit tool calls can enter the same typed
execution and state path, and the controller can always abstain.}
\label{fig:architecture}
\end{figure}

\subsection{Semantic interrupts are control events}

A \emph{semantic interrupt} is an automatically detected opportunity for
specialized assistance during generation.  It is not necessarily a hardware
interrupt, a token emitted by the model, or a pause after every sentence.  The
implementation may watch completed grammar fragments, parse text, inspect a
frozen hidden state, or eventually be trained jointly.  Three properties are
essential:
\begin{enumerate}
  \item the event is grounded in the developing inference state rather than
  only the original query;
  \item recognition is separated from typed normalization and execution; and
  \item every intervention, abstention, state update, and failure is observable.
\end{enumerate}

Explicit tool invocation remains available.  Heavyweight, risky, ambiguous, or
user-visible actions should generally require an explicit call or confirmation.
Automatic interrupts are most plausible for bounded, reversible, low-latency
operations with high-confidence semantics.

\subsection{Typed micro-IR}

Natural language should not be passed directly to an unrestricted executor when
a narrower interface is available.  Example computational events include
\[
\begin{split}
&\texttt{EVAL\_ARITH}(\texttt{op=MUL},\texttt{lhs=37},\texttt{rhs=48}),\\
&\texttt{GRAPH\_QUERY}(\texttt{op=REACHABLE},\texttt{src=A},\texttt{dst=D}),\\
&\texttt{CHECK\_SAT}(\texttt{constraints=C},\texttt{timeout=100ms}).
\end{split}
\]
Example epistemic events include
\[
\begin{split}
\texttt{RETRIEVE}(&\texttt{role=VERIFY\_CLAIM},
\texttt{entity=e},\texttt{relation=r},\\
&\texttt{source=DB},\texttt{as\_of=2026\mbox{-}08\mbox{-}28}).
\end{split}
\]
The schema should express uncertainty and refusal.  A normalizer may return
\textsc{ambiguous}; a router may return \textsc{none}; an executor may return
\textsc{timeout}, \textsc{unsupported}, or \textsc{no-evidence}.  This prevents
the runtime from turning weak interpretation into hard state.

\subsection{Micro-state and epistemic status}

Micro-state is the structured working record produced by interventions.  A
minimal record contains a stable identifier, type, payload, status, confidence,
source or engine version, timestamp, scope, parent dependencies, and the event
that created it.  Epistemic records should distinguish at least
\textsc{supported}, \textsc{contradicted}, \textsc{unverified}, \textsc{stale},
\textsc{ambiguous}, and \textsc{conflict}.  Computational records should retain
the formalization and execution status, not just the final value.

State is initially append-only and textually reinjected.  Retraction,
transactions, branch-local hypotheses, and selective native interfaces are
later hypotheses.  Introducing them before persistent-state failures are
observed would confound the early experiments.

\section{Coupling Ladder}

The architecture is a ladder of testable couplings rather than a requirement to
build the final system at once:
\begin{enumerate}
  \item \textbf{Strict syntax or grammar.} Completed arithmetic or controlled
  retrieval forms trigger deterministic handlers.
  \item \textbf{Lexical and syntactic rules.} Transparent patterns recognize a
  wider set of generation prefixes.
  \item \textbf{Heuristic routing.} Typed cues select among several engines or
  evidence sources.
  \item \textbf{Learned text-level interrupts.} Small detectors and parsers
  predict whether, what, and where to offload, with abstention and calibration.
  \item \textbf{Hidden-state routing.} Probes use frozen Transformer states when
  text alone is insufficient.
  \item \textbf{Joint interfaces.} The model and runtime co-adapt, potentially
  using structured or native key/value access.
\end{enumerate}

Each rung must beat the strongest simpler rung under matched capability and cost
or solve a measured failure that the simpler rung cannot.  ``Tighter'' is not a
synonym for ``better'': additional coupling can increase latency, false
interventions, opacity, and distribution-shift sensitivity.

\section{Illustrative Traces}

These examples define intended semantics; they are not implementation results.

\paragraph{Computational reflex.}
While producing ``The total is $37\times48=$'', a strict watcher recognizes a
completed bounded expression.  The normalizer emits
\texttt{EVAL\_ARITH(MUL,37,48)}, the calculator returns $1776$, and the state
records the expression, result, engine version, and execution status.  The
materializer exposes ``calculator result: 1776'' before decoding resumes.  A
quoted expression, malformed expression, or expression outside the budget must
not execute.

\paragraph{Epistemic reflex.}
While producing ``As of August 2026, the current policy limit is'', a watcher
marks an \textsc{assertion-pending} role with a freshness requirement.  A source
router chooses the versioned policy database; the source returns two conflicting
records.  The state is \textsc{conflict}, not \textsc{supported}, and the model
must cite the disagreement, ask for a governing jurisdiction, or abstain.

\paragraph{Composition.}
For ``Which region grew fastest?'', retrieval first obtains versioned regional
values.  A computational event then calculates growth rates whose dependencies
point to those records.  If a source record is retracted, downstream values are
marked invalid rather than left as apparently exact facts.  This dependency
behavior belongs to the transactional-state stage, not the initial reflex
experiments.

\section{Relationship to Prior Work}

The proposal is deliberately adjacent to several mature lines of work.  The
claim is architectural organization and a falsifiable program across them, not
priority over their constituent mechanisms.

\subsection{Tool use and agentic control}

Toolformer trains a model to decide which API to call, with what arguments, and
how to incorporate the result \cite{toolformer}.  ReAct interleaves model-written
reasoning and actions against external environments \cite{react}.  These works
establish that tools can improve model behavior and that action can occur within
an inference trajectory.  Our controller differs mainly in where responsibility
is placed: early semantic interrupts are runtime-recognized events with typed
state, while model-authored calls remain a fallback and a baseline.  This is an
empirical comparison, not an assumption that runtime control is superior.

\subsection{Program-aided and neuro-symbolic reasoning}

PAL delegates generated programs to an interpreter \cite{pal}; Logic-LM maps
problems to symbolic formulations and invokes deterministic solvers
\cite{logiclm}.  AlphaGeometry combines a neural language model with a symbolic
deduction engine \cite{alphageometry}, and LeanDojo studies language-model
interaction with the Lean theorem prover \cite{leandojo}.  Dynamic solver
composition routes natural-language problems among multiple logical paradigms
\cite{dynamic}.  Forethought composes neural and symbolic primitives into
inspectable programs \cite{forethought}.  These systems motivate typed,
verifiable external computation.  The proposed program focuses on incremental
generation-time detection, a common state contract spanning computation and
retrieval, and controlled progression from reflexes to learned coupling.

\subsection{Retrieval-augmented generation}

RAG combines parametric generation with non-parametric retrieval
\cite{rag}.  More directly, FLARE actively retrieves during long-form generation
by predicting an upcoming sentence, using low-confidence tokens to trigger
retrieval, querying with the prediction, and regenerating with evidence
\cite{flare}.  Self-RAG trains the generative LM to retrieve on demand and to
emit reflection tokens concerning retrieval relevance, generation support, and
response utility \cite{selfrag}.  Dynamic and adaptive retrieval methods further
study real-time information needs and question-complexity policies
\cite{dragin,adaptiverag}; corrective RAG evaluates retrieved evidence and can
fall back to web search \cite{crag}; multi-source adaptive RAG already studies
when, what, and which source to retrieve \cite{mspr}.  Generation-time or
learned adaptive retrieval therefore is not novel by itself.

The retrieval track instead asks whether confidence and semantic epistemic risk
fail on different cases, whether structurally different evidence substrates
require source-specific routing and query semantics, where learned control
should live, and whether typed persistent evidence state enables later
dependency propagation and retraction.

\begin{table}[p]
\centering
\scriptsize
\setlength{\tabcolsep}{2.2pt}
\caption{Required retrieval comparison. Entries describe the central published
mechanisms and the proposed program requirements, not claims that every system
is restricted to one deployment configuration.}
\label{tab:retrieval-comparison}
\begin{tabularx}{\linewidth}{@{}p{0.18\linewidth}YYYY@{}}
\toprule
Dimension & Conventional RAG & FLARE & Self-RAG & Cognitive-coprocessor program \\
\midrule
Retrieval timing & Mostly pre-generation & During generation & Adaptive during generation & During generation; upfront retrieval retained \\
Trigger & Request or harness & Low-confidence predicted tokens & Learned retrieval\slash reflection behavior & Semantic\slash epistemic event; later learned \\
Forward-looking state & Usually no & Predicted upcoming sentence & Adaptive segment-level generation & Prefix/text/semantic state; later latent \\
Learned trigger & Usually no & No & Yes & Paper 3.5 \\
Multiple retrieval engines & Orchestration-specific & Not central & Not central & Central in Papers 2.5/3.5 \\
Learned source routing & Not central & No & Not central & Paper 3.5 \\
DB\slash vector\slash web distinction & Orchestration-specific & No & No & First-class action semantics \\
Evidence\slash reflection state & Passages & Retrieved evidence & Relevance\slash support\slash utility reflection & Typed persistent epistemic micro-state \\
Calculator\slash logic\slash SMT & Separate tools & No & No & Parallel compute track \\
Automatic symbolic inference & No & No & No & Paper 2 onward \\
Transaction\slash retraction & No & No & Limited critique & Paper 4 \\
PRA\slash KV micro-state & No & No & No & Conditional Paper 5 interface \\
Target architecture & Retrieval augmentation & Active retrieval & Self-reflective adaptive RAG & Heterogeneous cognitive runtime \\
\bottomrule
\end{tabularx}
\end{table}

\subsection{Context and runtime management}

MemGPT uses operating-system ideas, memory tiers, and interrupts to virtualize
context \cite{memgpt}.  INFERCEPT addresses scheduling and key/value-cache costs
for interrupted augmented inference \cite{infercept}.  These works make clear
that interruption and persistent state are also runtime concerns.  Progressive
Retrieval Attention (PRA), developed as a separate research line in this
project, is a later candidate for selectively materializing large typed state;
it is neither assumed nor required in Papers 1--3.5.

\begin{table}[t]
\centering
\footnotesize
\setlength{\tabcolsep}{3pt}
\caption{Related-work matrix.  ``Model'' means the language model explicitly
authors or predicts the action; ``runtime'' means an external policy can
intervene.  Entries summarize the primary mechanism, not every possible
extension.}
\label{tab:related}
\begin{tabularx}{\linewidth}{@{}p{0.16\linewidth}p{0.14\linewidth}p{0.18\linewidth}YY@{}}
\toprule
Paradigm & Trigger/control & External operation & State returned & Main contrast here \\
\midrule
Function/tool calling & Model or orchestrator & Arbitrary APIs & Usually serialized observation & Automatic typed reflex is optional, not required \\
ReAct / planner--executor & Model-authored plan/action & Tools and environments & Text trajectory & Runtime can recognize bounded needs without a full plan \\
PAL / program aided & Prompted program generation & Interpreter & Program output & Incremental events and persistent typed state \\
Neuro-symbolic / theorem proving & Formalization plus solver policy & Logic, proof, geometry & Proof state or solver output & Twin compute/retrieval contract and staged interrupts \\
Upfront RAG & Input-time query & Usually one retrieval stack & Retrieved context & Needs can emerge after generation begins \\
FLARE / Self-RAG & Confidence or learned tokens & On-demand retrieval & Passages and reflection signals & Multiple source types plus shared compute/retrieval state \\
Memory virtualization & Model/runtime memory policy & Store and retrieve context & Memory pages or context & Operation/evidence semantics and provenance, not only capacity \\
Proposed program & Runtime first; learned later & Heterogeneous compute and retrieval & Typed, scoped, provenance-bearing micro-state & Joint abstraction tested rung by rung \\
\bottomrule
\end{tabularx}
\end{table}

\section{Novelty Claim and Non-Claims}

The strongest defensible claim is a conjunction.  No single row in
Table~\ref{tab:novelty} is asserted to be unprecedented.  The research question
is whether the conjunction creates measurable value and whether it can be built
incrementally without importing all later mechanisms into the first result.

\begin{table}[t]
\centering
\scriptsize
\caption{Novelty comparison.  The proposed entries are program requirements,
not completed empirical achievements.}
\label{tab:novelty}
\begin{tabularx}{\linewidth}{@{}p{0.23\linewidth}YYYY@{}}
\toprule
Dimension & Tool\slash agent systems & Active retrieval & Neuro-symbolic systems & Proposed program \\
\midrule
Separate competence from knowledge need & Usually implicit & Knowledge focus & Competence focus & Explicit twin taxonomy \\
Common typed event contract & API\slash action schema & Retrieval\slash reflection schema & Formal problem schema & Shared envelope, family-specific payloads \\
Generation-time automatic trigger & Sometimes model-authored & Central mechanism & Usually task- or plan-level & Runtime reflex first, learned later \\
Heterogeneous engine\slash source routing & Common in agents & Sometimes & Emerging & Both tracks, separately evaluated \\
Persistent provenance-bearing state & Often text history & Evidence passages\slash citations & Proof or solver state & Unified micro-state with status and dependencies \\
Retraction and branch scope & Agent-dependent & Rarely central & Solver-dependent & Explicit later evidence gate \\
Complexity staged by falsification & Not generally a series constraint & Not generally & Not generally & Required before each coupling rung \\
\bottomrule
\end{tabularx}
\end{table}

This paper does \emph{not} claim that symbolic systems replace Transformers,
that retrieved evidence is true, that formal derivations validate their
premises, that hidden-state interfaces are necessarily better than text, or that
the complete runtime has been implemented.  It also does not claim that the
term \emph{semantic interrupt} supersedes prior notions of active retrieval,
tool-use tokens, or runtime interruption; it names the common control event used
by this program.  In particular, neither retrieval during generation nor learned
adaptive retrieval is claimed as novel.  FLARE and Self-RAG are direct prior art.
Nor are relevance/support labels a Paper 4 contribution; Paper 4 is reserved for
persistent cross-source provenance, dependencies, transactional scopes,
retraction, and backtracking.

\section{Falsifiable Program Hypotheses}

Let $q$ denote task quality, $u$ unsupported-claim rate, $n$ generated tokens,
$m$ model calls, $\ell$ latency, $a$ external calls, and $s$ active state size.
For a system $j$, define a declared cost-sensitive utility
\[
  U_j=q_j-\lambda_u u_j-\lambda_n n_j-\lambda_m m_j
      -\lambda_\ell \ell_j-\lambda_a a_j-\lambda_s s_j,
\]
where all $\lambda$ values are fixed before evaluation and raw metrics are
reported alongside $U$.  No single weighting should decide the result; report
Pareto frontiers and threshold sweeps.

\paragraph{H1: Reliability.}
At matched model, task information, and external capabilities, accepted
computational interrupts reduce operation errors and accepted epistemic
interrupts reduce unsupported factual commitments without an offsetting rise in
false interventions.  H1 is falsified if component-correct interventions do not
improve end-task reliability, or if detection, normalization, routing, or
evidence-use errors erase the engine advantage.

\paragraph{H2: Efficiency.}
For regimes containing repeated bounded operations or emerging information
needs, automatic typed assistance yields a better quality--cost frontier than
neural emulation, verbose explicit calls, or indiscriminate retrieval.  H2 is
falsified if gains vanish under matched token, latency, call, and engine-work
budgets, or if a simpler baseline weakly dominates the proposed mechanism.

\paragraph{H3: Scaling.}
As reasoning depth $d$, problem size $p$, source count $r$, factual subgoals $g$,
or state size $s$ increases, a suitable heterogeneous system degrades more
slowly than matched homogeneous and explicit-tool baselines.  This requires
estimating slopes or fitted scaling curves, not comparing one endpoint.  H3 is
falsified if the proposed system's error or cost slope is no better, or if
routing and state-management overhead grows faster than the saved work.

\paragraph{H4: Small-model leverage.}
Smaller Transformers receive a larger marginal benefit from correct
coprocessor assistance because they can concentrate on recognizing structure
and using results rather than emulating every operation.  Test the interaction
between model size and intervention condition with matched model families and
tasks.  H4 is falsified if relative gains are flat, favor only larger models, or
depend on oracle triggers unavailable to the smaller models.

\paragraph{H5: Typed-state leverage.}
Persistent typed state improves multi-step use, auditability, and correction
once workloads require reuse or revision.  H5 is falsified if independent text
tool calls match quality and cost, if state is rarely reused, or if stale-state
contamination outweighs reuse benefits.  Transactional and native interfaces
are justified only after this simpler state hypothesis survives.

\section{Evaluation Protocol}

\subsection{Matched baselines}

Every empirical paper should include, where applicable: (i) the base model
without assistance; (ii) matched prompting; (iii) conventional explicit
tool/function calling; (iv) upfront RAG for knowledge tasks; (v) the proposed
mechanism; and (vi) oracle trigger, route, normalization, or result-use controls.
Strong planner--executor or active-retrieval baselines should be added when the
task reaches their intended regime.  Baselines must receive the same underlying
engine or corpus whenever the scientific question concerns interface or timing.
Paper 1.5 must include a FLARE-like confidence condition and confidence/semantic
OR and AND ablations.  Paper 2.5 must compare universal retrieval, heterogeneous
routing, explicit source choice, matched-cost broadcast, and oracle source
routing.  Paper 3.5 must compare FLARE-like confidence control, Self-RAG-style
main-model control, a text sidecar, a matched frozen-LM hidden-state probe, and
oracle control.  Approximations must be labeled rather than named as canonical
systems.

\subsection{Component decomposition}

An end-to-end score cannot identify whether a coprocessor architecture failed.
Each accepted and rejected candidate should be logged through
\[
\textsc{detect}\rightarrow\textsc{normalize}\rightarrow\textsc{route}
\rightarrow\textsc{execute/retrieve}\rightarrow\textsc{update}
\rightarrow\textsc{materialize}\rightarrow\textsc{use}.
\]
Report trigger precision/recall and false-intervention rate; normalization exact
or semantic correctness; route top-$1$/top-$k$ accuracy and abstention;
engine/retrieval correctness; evidence status and provenance correctness;
state-update correctness; result-use rate; and final-answer quality.  Oracle
controls estimate the headroom at each boundary.

\subsection{Cost, scaling, and robustness}

Report generated and reinjected tokens, model calls, intervention count, source
fanout, engine work, wall-clock and engine latency, active and total state size,
and monetary cost where meaningful.  Sweep problem depth/size, source count,
number of factual subgoals, context/state size, and model size.  Include hard
non-trigger controls, quoted and hypothetical statements, malformed events,
ambiguous requests, stale and conflicting evidence, unavailable engines,
timeouts, adversarial source text, and incorrect coprocessor outputs.

\subsection{Reproducibility and analysis}

Pin model, tokenizer, corpus/index, engine, and environment revisions.  Preserve
raw event traces and machine-readable summaries.  Use identity-disjoint splits
where templates or entities can leak, multiple seeds for stochastic headline
results, confidence intervals, and paired tests where examples are shared.
Predeclare primary outcomes and cost weights.  Qualitative traces illustrate
failure modes but are not evidence of success.

\section{Interleaved Research Roadmap}

The computational and epistemic tracks alternate so that each begins with the
smallest transparent reflex, expands to heterogeneity, and only then learns
interrupts.  Papers 4--7 integrate mechanisms only when earlier results expose a
need.

\begin{table}[p]
\centering
\scriptsize
\setlength{\tabcolsep}{2.5pt}
\caption{Per-paper hypothesis, falsification condition, and gate.}
\label{tab:roadmap}
\begin{tabularx}{\linewidth}{@{}p{0.08\linewidth}p{0.16\linewidth}YYY@{}}
\toprule
Paper & Mechanism & Primary hypothesis & Falsification / negative result & Evidence gate \\
\midrule
0 & Position and protocol & Twin typed coprocessor tracks form a coherent testable program & No discriminating baselines, metrics, or experiments follow from the distinction & Proceed only with scoped empirical tests \\
1 & Strict reflex calculator & Automatic bounded arithmetic improves reliability or cost over model-only and explicit calls & Explicit calls match the full frontier, or false triggers erase exact execution gains & Heterogeneous compute needs must be observed before Paper 2 \\
1.5 & Confidence versus semantic one-source retrieval & Epistemic-risk rules catch evidence needs missed by FLARE-like uncertainty, especially confident commitments & Confidence catches the same cases at equal/lower cost; semantic rules duplicate uncertainty & Multiple sources must address a measured limitation before Paper 2.5 \\
2 & Multiple computational engines and typed state & Typed routing/state beats a single engine or repeated explicit calls on mixed tasks & Router overhead/errors dominate; state is not reused & Brittle recognition must limit performance before Paper 3 \\
2.5 & Heterogeneous evidence substrates, heuristic routing & Source-specific semantics improve relevance/cost over universal retrieval or broadcast & Oracle routing adds little; universal retrieval or matched-cost broadcast dominates & Oracle heterogeneous routing must win before Paper 3.5 \\
3 & Learned computational interrupts & Learned detection/normalization generalizes beyond syntax with calibrated abstention & Heuristics match OOD performance/cost; learned false interrupts dominate & Persistent compute state must create revision pressure before Paper 4 \\
3.5 & Main-model versus sidecar epistemic control & Learned control routes heterogeneous sources; transfer/cost reveal where competence should live & Self-RAG-style control dominates, sidecars do not transfer, or source routing adds no value & Both tracks must exhibit concrete provenance/revision failures \\
4 & Transactional unified state & Persistent provenance, dependencies, scope, retraction, and rollback reduce contamination & Append-only state matches quality; rollback errors exceed contamination avoided & Text state must become a measured access bottleneck \\
5 & Structured, PRA, and native interfaces & Selective materialization improves quality--memory--latency tradeoffs at large state & Text serialization remains competitive; native access adds no matched benefit & Co-adaptation requires stable useful interfaces first \\
6 & Joint learning to offload and acquire & Co-adaptation improves formulation/result use and reduces redundant neural work & Gains vanish without coprocessors or under robustness tests; no division-of-labor change is measurable & Integrate only mechanisms with replicated gains \\
7 & Heterogeneous runtime & Validated components compose on mixed workloads with graceful fallback & Integration is dominated by a strong agent/tool baseline or interactions negate component gains & Final synthesis; unsupported components remain excluded \\
\bottomrule
\end{tabularx}
\end{table}

\section{Failure Modes, Safety, and Governance}

Automatic intervention creates risks beyond ordinary answer error.
\begin{itemize}
  \item \textbf{False interrupts and missed needs.} Over-triggering wastes work
  and can alter correct generations; under-triggering preserves the original
  weakness.  Threshold and selective-risk curves are therefore mandatory.
  \item \textbf{Semantic misnormalization.} A correct engine can make an
  incorrect formalization look authoritative.  Store the source span and typed
  event, expose uncertainty, and use oracle-normalization controls.
  \item \textbf{Unsafe execution.} Prefer bounded domain-specific engines over
  arbitrary generated code.  Enforce time, memory, recursion, expression-size,
  network, and side-effect budgets.
  \item \textbf{Evidence laundering.} Retrieval relevance is not source
  reliability, and exact computation over a retrieved value does not strengthen
  that value's provenance.  Derived records must retain source dependencies.
  \item \textbf{Prompt injection and source poisoning.} Evidence is data, not
  authority to modify runtime policy.  Separate source content from control
  channels and test adversarial documents.
  \item \textbf{Stale or branch-leaked state.} Hypotheticals, quotations, and
  abandoned plans must not silently become global facts.  Early papers constrain
  such cases; Paper 4 tests scope and rollback directly.
  \item \textbf{Latency and denial of service.} Interrupt storms and expensive
  routing can dominate generation.  Apply per-event and per-request budgets,
  deduplication, caching, circuit breakers, and explicit fallback behavior.
\end{itemize}

Users should be able to inspect when external systems were consulted, what
state was created, and which claims depend on it.  High-impact external actions
remain outside automatic reflexes unless separately authorized; the present
program concerns computation and evidence acquisition, not autonomous execution
of consequential real-world actions.

\section{Limitations and Open Questions}

First, the competence/knowledge distinction is operational rather than
ontological, and composed tasks can blur it.  Its value must be judged by error
localization and experimental predictions.  Second, observing a generation
prefix may be too late to prevent a commitment or too early to formulate the
right query.  Prospective interruption and retrospective verification should be
compared.  Third, pauses can perturb decoding independently of the returned
result; sham-interrupt controls can estimate this effect.  Fourth, a shared
event envelope may simplify observability while family-specific semantics still
require distinct normalizers, engines, and trust policies.

Fifth, persistent state can become a second context-management problem.  Text
reinjection is the correct initial baseline, while structured selection,
Progressive Retrieval Attention, and native key/value interfaces remain
conditional Paper 5 hypotheses.  Sixth, learned interrupts may inherit model
biases and fail under domain shift; calibration, abstention, and explicit tools
must remain available.  Finally, the architecture may lose to a well-engineered
agent with ordinary calls.  That outcome would be scientifically useful: it
would identify where tighter coupling adds complexity without value.

\section{Conclusion}

The proposed research program treats a language model neither as a complete
reasoning machine nor as a thin dispatcher.  The Transformer remains the central
learned substrate for interpreting language, recognizing structure, composing
results, and communicating.  Specialized systems contribute operations and
evidence under explicit typed contracts.  Separating competence from knowledge
deficits yields twin computational and epistemic tracks; semantic interrupts,
micro-IR, and provenance-bearing state provide a common experimental interface.

The proposal succeeds only if this organization produces reproducible gains in
reliability, cost, scaling, or small-model leverage against strong simpler
baselines.  The staged roadmap is therefore part of the architecture: begin
with strict reflexes, measure their limits, and add routing, learning,
transactions, native access, and co-adaptation only when evidence requires them.

\begin{thebibliography}{99}

\bibitem{toolformer}
T.~Schick, J.~Dwivedi-Yu, R.~Dessi, R.~Raileanu, M.~Lomeli, L.~Zettlemoyer,
N.~Cancedda, and T.~Scialom.
Toolformer: Language models can teach themselves to use tools.
\emph{NeurIPS}, 2023.
\url{https://arxiv.org/abs/2302.04761}.

\bibitem{react}
S.~Yao, J.~Zhao, D.~Yu, N.~Du, I.~Shafran, K.~Narasimhan, and Y.~Cao.
ReAct: Synergizing reasoning and acting in language models.
\emph{ICLR}, 2023.
\url{https://arxiv.org/abs/2210.03629}.

\bibitem{pal}
L.~Gao, A.~Madaan, S.~Zhou, U.~Alon, P.~Liu, Y.~Yang, J.~Callan, and G.~Neubig.
PAL: Program-aided language models.
\emph{ICML}, 2023.
\url{https://arxiv.org/abs/2211.10435}.

\bibitem{rag}
P.~Lewis, E.~Perez, A.~Piktus, F.~Petroni, V.~Karpukhin, N.~Goyal,
H.~Kuttler, M.~Lewis, W.-t.~Yih, T.~Rocktaschel, S.~Riedel, and D.~Kiela.
Retrieval-augmented generation for knowledge-intensive NLP tasks.
\emph{NeurIPS}, 2020.
\url{https://arxiv.org/abs/2005.11401}.

\bibitem{flare}
Z.~Jiang, F.~F.~Xu, L.~Gao, Z.~Sun, Q.~Liu, J.~Dwivedi-Yu, Y.~Yang,
J.~Callan, and G.~Neubig.
Active retrieval augmented generation.
\emph{EMNLP}, 2023.
\url{https://arxiv.org/abs/2305.06983}.

\bibitem{selfrag}
A.~Asai, Z.~Wu, Y.~Wang, A.~Sil, and H.~Hajishirzi.
Self-RAG: Learning to retrieve, generate, and critique through self-reflection.
\emph{ICLR}, 2024.
\url{https://arxiv.org/abs/2310.11511}.

\bibitem{dragin}
W.~Su, Y.~Tang, Q.~Ai, Z.~Wu, and Y.~Liu.
DRAGIN: Dynamic retrieval augmented generation based on the information needs of
large language models.
arXiv:2403.10081, 2024.
\url{https://arxiv.org/abs/2403.10081}.

\bibitem{adaptiverag}
S.~Jeong, J.~Baek, S.~Cho, S.~J.~Hwang, and J.~C.~Park.
Adaptive-RAG: Learning to adapt retrieval-augmented large language models
through question complexity.
arXiv:2403.14403, 2024.
\url{https://arxiv.org/abs/2403.14403}.

\bibitem{crag}
S.-Q.~Yan, J.-C.~Gu, Y.~Zhu, and Z.-H.~Ling.
Corrective retrieval augmented generation.
arXiv:2401.15884, 2024.
\url{https://arxiv.org/abs/2401.15884}.

\bibitem{mspr}
Q.~Zhao, R.~Wang, X.~Wang, D.~Zha, and N.~Mu.
Towards multi-source retrieval-augmented generation via synergizing reasoning
and preference-driven retrieval.
arXiv:2411.00689, 2024.
\url{https://arxiv.org/abs/2411.00689}.

\bibitem{logiclm}
L.~Pan, A.~Albalak, X.~Wang, and W.~Y.~Wang.
Logic-LM: Empowering large language models with symbolic solvers for faithful
logical reasoning.
\emph{Findings of EMNLP}, 2023.
\url{https://arxiv.org/abs/2305.12295}.

\bibitem{alphageometry}
T.~H.~Trinh, Y.~Wu, Q.~V.~Le, H.~He, and T.~Luong.
Solving olympiad geometry without human demonstrations.
\emph{Nature}, 2024.
\url{https://arxiv.org/abs/2312.11871}.

\bibitem{leandojo}
K.~Yang, A.~M.~Swope, A.~Gu, R.~Chalamala, P.~Song, S.~Yu, S.~Godil,
R.~J.~Prenger, and A.~Anandkumar.
LeanDojo: Theorem proving with retrieval-augmented language models.
\emph{NeurIPS}, 2023.
\url{https://arxiv.org/abs/2306.15626}.

\bibitem{dynamic}
L.~Xu, P.~Beckmann, M.~Valentino, and A.~Freitas.
Adaptive LLM-symbolic reasoning via dynamic logical solver composition.
arXiv:2510.06774, 2025.
\url{https://arxiv.org/abs/2510.06774}.

\bibitem{forethought}
V.~Bhat, J.~Vaghasiya, and E.~A.~Gonzalez.
Forethought: Verifiable reasoning from neurosymbolic primitive programming.
arXiv:2607.04096, 2026.
\url{https://arxiv.org/abs/2607.04096}.

\bibitem{memgpt}
C.~Packer, S.~Wooders, K.~Lin, V.~Fang, S.~G.~Patil, I.~Stoica, and
J.~E.~Gonzalez.
MemGPT: Towards LLMs as operating systems.
arXiv:2310.08560, 2023.
\url{https://arxiv.org/abs/2310.08560}.

\bibitem{infercept}
R.~Abhyankar, Z.~He, V.~Srivatsa, H.~Zhang, and Y.~Zhang.
InferCept: Efficient intercept support for augmented large language model
inference.
arXiv:2402.01869, 2024.
\url{https://arxiv.org/abs/2402.01869}.

\end{thebibliography}

\end{document}
